\documentclass{beamer}

\usepackage{beamerthemeSingapore}
\usepackage[brazil]{babel}
\usepackage[T1]{fontenc}
\usepackage[utf8]{inputenc}
\usepackage{graphicx}
\usepackage{url}
\usepackage{xcolor}
\usepackage{booktabs}
\usepackage{array}
\usepackage{ragged2e}
\usepackage{etoolbox}
\usepackage{amsmath}
\usepackage{minted}
\usepackage{tikz}
\usetikzlibrary{positioning, arrows.meta, calc}
% Grafos compartilhados com o roteiro.
% Grafos da aula de representação e percursos. Compartilhado entre o roteiro
% (figuras.tex) e os slides (grafos.tex da pasta 00_grafos). Os grafos são os
% mesmos dos programas da pasta codes.

\tikzset{
  gv/.style = {circle, draw=black, fill=white, inner sep=0pt,
    minimum size=1.6em, line width=0.6pt, font=\sffamily\small},
  ga/.style = {draw=black!55, line width=0.7pt},
  garv/.style = {draw=blue!70!black, line width=2.2pt},
  gnota/.style = {font=\sffamily\scriptsize, text=blue!70!black},
  arco/.style = {-{Stealth[length=2.4mm]}, draw=black!70, line width=0.7pt},
  pv/.style = {rectangle, rounded corners=2pt, draw=black, fill=white,
    inner sep=3pt, line width=0.6pt, font=\ttfamily\small},
  arcov/.style = {-{Stealth[length=2.4mm]}, draw=red!75!black, line width=1.6pt}
}

% ------------------------------------------------------------------
% Grafo G, de bfs.go, dfs.go e representacao.go.
\newcommand{\Gnos}{%
  \node [gv] (g0) at (0, 1.5) {0};
  \node [gv] (g1) at (1.5, 3) {1};
  \node [gv] (g2) at (1.5, 1.5) {2};
  \node [gv] (g5) at (1.5, 0) {5};
  \node [gv] (g3) at (3, 3) {3};
  \node [gv] (g4) at (3, 0) {4};
  \node [gv] (g6) at (4.5, 1.5) {6};
  \node [gv] (g7) at (6, 1.5) {7};
}
\newcommand{\Garestas}{%
  \foreach \a/\b in {0/1, 0/2, 0/5, 1/3, 2/3, 2/4, 3/6, 4/5, 4/6, 6/7}
    \draw [ga] (g\a) -- (g\b);
}

\newcommand{\grafoG}{%
\begin{tikzpicture}
  \Gnos
  \Garestas
\end{tikzpicture}}

% Árvore da busca em largura a partir de 0, com a distância de cada vértice.
\newcommand{\grafoGbfs}{%
\begin{tikzpicture}
  \Gnos
  \Garestas
  \foreach \a/\b in {0/1, 0/2, 0/5, 1/3, 2/4, 3/6, 6/7}
    \draw [garv] (g\a) -- (g\b);
  \Gnos
  \foreach \v/\d/\p in {0/0/left, 1/1/above, 2/1/right, 5/1/below,
                        3/2/above, 4/2/below, 6/3/below, 7/4/above}
    \node [gnota, \p=1pt of g\v] {d=\d};
\end{tikzpicture}}

% Árvore da busca em profundidade a partir de 0, com a ordem de visita.
\newcommand{\grafoGdfs}{%
\begin{tikzpicture}
  \Gnos
  \Garestas
  \foreach \a/\b in {0/1, 1/3, 3/2, 2/4, 4/5, 4/6, 6/7}
    \draw [garv] (g\a) -- (g\b);
  \Gnos
  \foreach \v/\d/\p in {0/1/left, 1/2/above, 3/3/above, 2/4/right,
                        4/5/below, 5/6/below, 6/7/below, 7/8/above}
    \node [gnota, \p=1pt of g\v] {\d\textordmasculine};
\end{tikzpicture}}

% ------------------------------------------------------------------
% Grafo de componentes.go: 10 vértices, 4 componentes.
\newcommand{\grafoComp}{%
\begin{tikzpicture}
  \node [gv] (h0) at (0, 1.6) {0};
  \node [gv] (h1) at (-0.7, 0.4) {1};
  \node [gv] (h2) at (0.7, 0.4) {2};
  \node [gv] (h3) at (2.2, 1.6) {3};
  \node [gv] (h4) at (2.2, 0.4) {4};
  \node [gv] (h5) at (3.7, 1.6) {5};
  \node [gv] (h6) at (5.0, 1.6) {6};
  \node [gv] (h7) at (3.7, 0.4) {7};
  \node [gv] (h8) at (5.0, 0.4) {8};
  \node [gv] (h9) at (6.5, 1.0) {9};
  \foreach \a/\b in {0/1, 0/2, 1/2, 3/4, 5/6, 5/7, 6/8, 7/8}
    \draw [ga] (h\a) -- (h\b);
  \node [gnota, below=0.25cm of h1, xshift=0.7cm] {id 0};
  \node [gnota, below=0.25cm of h4] {id 1};
  \node [gnota, below=0.25cm of h7, xshift=0.65cm] {id 2};
  \node [gnota, below=0.85cm of h9] {id 3};
\end{tikzpicture}}

% ------------------------------------------------------------------
% Pacotes de topologica.go. O arco B -> A significa que A importa B.
\newcommand{\Pnos}{%
  \node [pv] (modelo) at (0, 2.4) {modelo};
  \node [pv] (log) at (0, 1.2) {log};
  \node [pv] (config) at (0, 0) {config};
  \node [pv] (banco) at (2.6, 1.2) {banco};
  \node [pv] (servico) at (5.2, 2.4) {servico};
  \node [pv] (http) at (7.8, 1.2) {http};
  \node [pv] (main) at (10.4, 1.2) {main};
}
\newcommand{\Parcos}{%
  \draw [arco] (modelo) -- (banco);
  \draw [arco] (log) -- (banco);
  \draw [arco] (config) -- (banco);
  \draw [arco] (banco) -- (servico);
  \draw [arco] (modelo) -- (servico);
  \draw [arco] (log) -- (servico);
  \draw [arco] (log) to [bend right=14] (http);
  \draw [arco] (servico) -- (http);
  \draw [arco] (servico) -- (main);
  \draw [arco] (http) -- (main);
  \draw [arco] (config) to [bend right=8] (main);
}
\newcommand{\grafoPacotes}{%
\begin{tikzpicture}
  \Pnos
  \Parcos
\end{tikzpicture}}

% O mesmo grafo com o arco servico -> modelo, que fecha um ciclo.
\newcommand{\grafoPacotesCiclo}{%
\begin{tikzpicture}
  \Pnos
  \Parcos
  \draw [arcov] (modelo) -- (banco);
  \draw [arcov] (banco) -- (servico);
  \draw [arcov] (servico) to [bend right=25] (modelo);
\end{tikzpicture}}

% A ordem topológica encontrada, da esquerda para a direita: todo arco aponta
% para a direita.
\newcommand{\grafoPacotesOrdem}{%
\begin{tikzpicture}
  \foreach \n/\x in {log/0, config/1.7, modelo/3.4, banco/5.1,
                     servico/6.8, http/8.5, main/10.2}
    \node [pv] (\n) at (\x, 0) {\n};
  \draw [arco] (modelo) -- (banco);
  \draw [arco] (banco) -- (servico);
  \draw [arco] (servico) -- (http);
  \draw [arco] (http) -- (main);
  \draw [arco] (log.north) .. controls +(0, 1.0) and +(0, 1.0) .. (banco.north);
  \draw [arco] (log.north) .. controls +(0, 1.4) and +(0, 1.4) .. (servico.north);
  \draw [arco] (log.north) .. controls +(0, 1.8) and +(0, 1.8) .. (http.north);
  \draw [arco] (servico.north) .. controls +(0, 0.8) and +(0, 0.8) .. (main.north);
  \draw [arco] (config.south) .. controls +(0, -0.8) and +(0, -0.8) .. (banco.south);
  \draw [arco] (modelo.south) .. controls +(0, -1.1) and +(0, -1.1) .. (servico.south);
  \draw [arco] (config.south) .. controls +(0, -1.6) and +(0, -1.6) .. (main.south);
\end{tikzpicture}}

% ------------------------------------------------------------------
% Componentes fortemente conexas: três componentes, destacadas pelo fundo.
\newcommand{\grafoCFC}{%
\begin{tikzpicture}
  \fill [orange!18, rounded corners=8pt] (-0.55, -0.6) rectangle (2.05, 2.1);
  \fill [green!18, rounded corners=8pt] (3.15, -0.6) rectangle (5.75, 2.1);
  \fill [blue!12, rounded corners=8pt] (6.85, 0.15) rectangle (7.95, 1.35);
  \node [gv] (c0) at (0, 1.5) {0};
  \node [gv] (c1) at (1.5, 1.5) {1};
  \node [gv] (c2) at (0.75, 0) {2};
  \node [gv] (c3) at (3.7, 1.5) {3};
  \node [gv] (c4) at (5.2, 1.5) {4};
  \node [gv] (c5) at (4.45, 0) {5};
  \node [gv] (c6) at (7.4, 0.75) {6};
  \draw [arco] (c0) -- (c1);
  \draw [arco] (c1) -- (c2);
  \draw [arco] (c2) -- (c0);
  \draw [arco] (c3) -- (c4);
  \draw [arco] (c4) -- (c5);
  \draw [arco] (c5) -- (c3);
  \draw [arco] (c1) -- (c3);
  \draw [arco] (c2) -- (c5);
  \draw [arco] (c4) -- (c6);
\end{tikzpicture}}

% ------------------------------------------------------------------
% Execução passo a passo das duas buscas em G, a partir de 0. Os estados
% seguem a saída de bfs.go e dfs.go. Cores: branco, ainda não alcançado;
% cinza, na fila (largura) ou com a chamada em andamento (profundidade); preto,
% já retirado da fila ou com a chamada terminada. A borda vermelha marca o
% vértice em processamento. Em azul, as arestas da árvore; em laranja, as
% arestas que entraram na árvore no passo.
\tikzset{
  gcinza/.style = {fill=black!25},
  gpreto/.style = {fill=black!80, text=white},
  gatual/.style = {draw=red!80!black, line width=2pt},
  gnova/.style = {draw=orange!90!black, line width=2.6pt},
  gcaixa/.style = {rectangle, draw=black, fill=black!25, line width=0.6pt,
    minimum width=0.55cm, minimum height=0.5cm, inner sep=0pt,
    font=\sffamily\small},
  gquadro/.style = {rectangle, draw=black, fill=black!25, line width=0.6pt,
    minimum width=1.7cm, minimum height=0.36cm, inner sep=0pt,
    font=\ttfamily\scriptsize},
  gtexto/.style = {anchor=west, font=\sffamily\footnotesize, inner sep=0pt},
  gtabela/.style = {anchor=east, font=\ttfamily\small, inner sep=0pt},
  glegenda/.style = {anchor=west, font=\sffamily\scriptsize, inner sep=0pt}
}

% Posição do rótulo de distância em volta de cada vértice, em graus.
\newcommand{\Gang}[1]{\ifcase#1 180\or 90\or 0\or 90\or 270\or 270\or 270\or 90\fi}

% Grafo G no estado de um passo. #1: vértice em processamento (pode ser
% vazio); #2: vértices cinza; #3: vértices pretos; #4: arestas da árvore;
% #5: arestas novas no passo.
\newcommand{\Gestado}[5]{%
  \Gnos
  \Garestas
  \foreach \a/\b in {#4} \draw [garv] (g\a) -- (g\b);
  \foreach \a/\b in {#5} \draw [gnova] (g\a) -- (g\b);
  \Gnos
  \foreach \v in {#2} \node [gv, gcinza] at (g\v) {\v};
  \foreach \v in {#3} \node [gv, gpreto] at (g\v) {\v};
  \foreach \v in {#1} \draw [gatual] (g\v) circle (0.8em);
}

% Um passo da busca em largura. #1 a #5 como em \Gestado; #6: fila, do início
% para o fim; #7: distâncias conhecidas, em pares v/d; #8: o que o passo faz.
\newcommand{\passoBFS}[8]{%
\begin{tikzpicture}
  \useasboundingbox (-1.1, -2.55) rectangle (8.8, 3.65);
  \Gestado{#1}{#2}{#3}{#4}{#5}
  \foreach \v/\d in {#7} {
    \edef\ang{\Gang{\v}}
    \path (g\v) ++(\ang:0.66 and 0.43) node [gnota] {d=\d};
  }
  \node [gtabela] at (7.0, 3.4) {v};
  \node [gtabela] at (7.95, 3.4) {dist};
  \node [gtabela] at (8.75, 3.4) {pai};
  \draw [black!60] (6.6, 3.18) -- (8.8, 3.18);
  \foreach \v in {0, ..., 7}
    \node [gtabela] at (7.0, 2.92 - 0.38*\v) {\v};
  \foreach \v/\d in {#7}
    \node [gtabela] at (7.95, 2.92 - 0.38*\v) {\d};
  \node [gtabela] at (8.75, 2.92) {-1};
  \foreach \a/\b in {#4} \node [gtabela] at (8.75, 2.92 - 0.38*\b) {\a};
  \foreach \a/\b in {#5} \node [gtabela] at (8.75, 2.92 - 0.38*\b) {\a};
  \node [gtexto] at (-1.1, -1.05) {fila:};
  \foreach \v [count=\i] in {#6}
    \node [gcaixa] at (-0.2 + 0.6*\i, -1.05) {\v};
  \node [gtexto] at (-1.1, -1.75) {#8};
  \node [gv, gcinza, minimum size=1em] at (-0.9, -2.35) {};
  \node [glegenda] at (-0.65, -2.35) {na fila};
  \node [gv, gpreto, minimum size=1em] at (0.9, -2.35) {};
  \node [glegenda] at (1.15, -2.35) {já retirado};
  \node [gv, gatual, minimum size=1em] at (3.2, -2.35) {};
  \node [glegenda] at (3.45, -2.35) {retirado agora};
\end{tikzpicture}}

% Um passo da busca em profundidade. #1 a #5 como em \Gestado, com #2 também
% sendo a pilha de chamadas, da base para o topo; #6: pré-ordem até o passo;
% #7: pós-ordem até o passo; #8: o que o passo faz.
\newcommand{\passoDFS}[8]{%
\begin{tikzpicture}
  \useasboundingbox (-1.1, -2.55) rectangle (8.8, 3.65);
  \Gestado{#1}{#2}{#3}{#4}{#5}
  \gdef\Gtopo{0}
  \foreach \v [count=\i] in {#2} {
    \node [gquadro] at (7.9, -0.2 + 0.42*\i) {visita(\v)};
    \xdef\Gtopo{\i}
  }
  \ifnum\Gtopo>0
    \node [gquadro, gatual, fill=none] at (7.9, -0.2 + 0.42*\Gtopo) {};
  \fi
  \node [glegenda, anchor=center] at (7.9, -0.25) {pilha de chamadas};
  \node [gtexto] at (-1.1, -0.85) {pré-ordem: \texttt{#6}};
  \node [gtexto] at (-1.1, -1.3) {pós-ordem: \texttt{#7}};
  \node [gtexto] at (-1.1, -1.75) {#8};
  \node [gv, gcinza, minimum size=1em] at (-0.9, -2.35) {};
  \node [glegenda] at (-0.65, -2.35) {chamada em andamento};
  \node [gv, gpreto, minimum size=1em] at (3.0, -2.35) {};
  \node [glegenda] at (3.25, -2.35) {chamada terminada};
  \node [gv, gatual, minimum size=1em] at (6.4, -2.35) {};
  \node [glegenda] at (6.65, -2.35) {em execução};
\end{tikzpicture}}

% Passos numerados a partir de 1, para \only nos slides.
\newcommand{\bfsPasso}[1]{\ifcase#1
\or\passoBFS{}{0}{}{}{}{0}{0/0}{início: marca o 0 e o põe na fila}
\or\passoBFS{0}{1, 2, 5}{}{}{0/1, 0/2, 0/5}{1, 2, 5}{0/0, 1/1, 2/1, 5/1}{retira 0, lista \texttt{[1 2 5]}: 1, 2 e 5 entram na fila}
\or\passoBFS{1}{2, 5, 3}{0}{0/1, 0/2, 0/5}{1/3}{2, 5, 3}{0/0, 1/1, 2/1, 3/2, 5/1}{retira 1, lista \texttt{[0 3]}: 3 entra na fila}
\or\passoBFS{2}{5, 3, 4}{0, 1}{0/1, 0/2, 0/5, 1/3}{2/4}{5, 3, 4}{0/0, 1/1, 2/1, 3/2, 4/2, 5/1}{retira 2, lista \texttt{[0 3 4]}: 4 entra na fila}
\or\passoBFS{5}{3, 4}{0, 1, 2}{0/1, 0/2, 0/5, 1/3, 2/4}{}{3, 4}{0/0, 1/1, 2/1, 3/2, 4/2, 5/1}{retira 5, lista \texttt{[0 4]}: todos marcados, nada entra}
\or\passoBFS{3}{4, 6}{0, 1, 2, 5}{0/1, 0/2, 0/5, 1/3, 2/4}{3/6}{4, 6}{0/0, 1/1, 2/1, 3/2, 4/2, 5/1, 6/3}{retira 3, lista \texttt{[1 2 6]}: 6 entra na fila}
\or\passoBFS{4}{6}{0, 1, 2, 5, 3}{0/1, 0/2, 0/5, 1/3, 2/4, 3/6}{}{6}{0/0, 1/1, 2/1, 3/2, 4/2, 5/1, 6/3}{retira 4, lista \texttt{[2 5 6]}: todos marcados, nada entra}
\or\passoBFS{6}{7}{0, 1, 2, 5, 3, 4}{0/1, 0/2, 0/5, 1/3, 2/4, 3/6}{6/7}{7}{0/0, 1/1, 2/1, 3/2, 4/2, 5/1, 6/3, 7/4}{retira 6, lista \texttt{[3 4 7]}: 7 entra na fila}
\or\passoBFS{7}{}{0, 1, 2, 5, 3, 4, 6}{0/1, 0/2, 0/5, 1/3, 2/4, 3/6, 6/7}{}{}{0/0, 1/1, 2/1, 3/2, 4/2, 5/1, 6/3, 7/4}{retira 7, lista \texttt{[6]}: fila vazia, fim da busca}
\fi}
\newcommand{\dfsPasso}[1]{\ifcase#1
\or\passoDFS{0}{0}{}{}{}{0}{}{\texttt{visita(0)}: lista \texttt{[1 2 5]}, desce para o 1}
\or\passoDFS{1}{0, 1}{}{}{0/1}{0 1}{}{\texttt{visita(1)}: lista \texttt{[0 3]}, 0 já marcado, desce para o 3}
\or\passoDFS{3}{0, 1, 3}{}{0/1}{1/3}{0 1 3}{}{\texttt{visita(3)}: lista \texttt{[1 2 6]}, 1 já marcado, desce para o 2}
\or\passoDFS{2}{0, 1, 3, 2}{}{0/1, 1/3}{3/2}{0 1 3 2}{}{\texttt{visita(2)}: lista \texttt{[0 3 4]}, 0 e 3 já marcados, desce para o 4}
\or\passoDFS{4}{0, 1, 3, 2, 4}{}{0/1, 1/3, 3/2}{2/4}{0 1 3 2 4}{}{\texttt{visita(4)}: lista \texttt{[2 5 6]}, 2 já marcado, desce para o 5}
\or\passoDFS{5}{0, 1, 3, 2, 4, 5}{}{0/1, 1/3, 3/2, 2/4}{4/5}{0 1 3 2 4 5}{}{\texttt{visita(5)}: lista \texttt{[0 4]}, todos marcados}
\or\passoDFS{4}{0, 1, 3, 2, 4}{5}{0/1, 1/3, 3/2, 2/4, 4/5}{}{0 1 3 2 4 5}{5}{\texttt{termina 5}: volta ao 4; o 6 não está marcado}
\or\passoDFS{6}{0, 1, 3, 2, 4, 6}{5}{0/1, 1/3, 3/2, 2/4, 4/5}{4/6}{0 1 3 2 4 5 6}{5}{\texttt{visita(6)}: lista \texttt{[3 4 7]}, 3 e 4 já marcados, desce para o 7}
\or\passoDFS{7}{0, 1, 3, 2, 4, 6, 7}{5}{0/1, 1/3, 3/2, 2/4, 4/5, 4/6}{6/7}{0 1 3 2 4 5 6 7}{5}{\texttt{visita(7)}: lista \texttt{[6]}, todos marcados}
\or\passoDFS{6}{0, 1, 3, 2, 4, 6}{5, 7}{0/1, 1/3, 3/2, 2/4, 4/5, 4/6, 6/7}{}{0 1 3 2 4 5 6 7}{5 7}{\texttt{termina 7}: volta ao 6; fim da lista}
\or\passoDFS{4}{0, 1, 3, 2, 4}{5, 7, 6}{0/1, 1/3, 3/2, 2/4, 4/5, 4/6, 6/7}{}{0 1 3 2 4 5 6 7}{5 7 6}{\texttt{termina 6}: volta ao 4; fim da lista}
\or\passoDFS{2}{0, 1, 3, 2}{5, 7, 6, 4}{0/1, 1/3, 3/2, 2/4, 4/5, 4/6, 6/7}{}{0 1 3 2 4 5 6 7}{5 7 6 4}{\texttt{termina 4}: volta ao 2; fim da lista}
\or\passoDFS{3}{0, 1, 3}{5, 7, 6, 4, 2}{0/1, 1/3, 3/2, 2/4, 4/5, 4/6, 6/7}{}{0 1 3 2 4 5 6 7}{5 7 6 4 2}{\texttt{termina 2}: volta ao 3; 6 já marcado, fim da lista}
\or\passoDFS{1}{0, 1}{5, 7, 6, 4, 2, 3}{0/1, 1/3, 3/2, 2/4, 4/5, 4/6, 6/7}{}{0 1 3 2 4 5 6 7}{5 7 6 4 2 3}{\texttt{termina 3}: volta ao 1; fim da lista}
\or\passoDFS{0}{0}{5, 7, 6, 4, 2, 3, 1}{0/1, 1/3, 3/2, 2/4, 4/5, 4/6, 6/7}{}{0 1 3 2 4 5 6 7}{5 7 6 4 2 3 1}{\texttt{termina 1}: volta ao 0; 2 e 5 já marcados, fim da lista}
\or\passoDFS{}{}{5, 7, 6, 4, 2, 3, 1, 0}{0/1, 1/3, 3/2, 2/4, 4/5, 4/6, 6/7}{}{0 1 3 2 4 5 6 7}{5 7 6 4 2 3 1 0}{\texttt{termina 0}: fim da busca}
\fi}

% ------------------------------------------------------------------
% Balde de tinta: imagem de 10 por 7 pixels, com 4 vizinhos por pixel. #1 = 0
% mostra a imagem antes do clique; #1 = 1, depois, com a componente do pixel
% clicado pintada.
\newcommand{\imagemBalde}[1]{%
\begin{tikzpicture}[x=0.38cm, y=0.38cm]
  \foreach \x/\y in {2/2, 3/2, 4/2, 5/2, 6/2, 7/2, 2/3, 3/3, 6/3, 7/3, 2/4, 3/4, 4/4, 5/4, 6/4, 7/4}
    {\ifnum#1=1 \fill [orange!70] (\x, -\y) rectangle ++(1, -1);\fi}
  \foreach \x/\y in {2/1, 3/1, 4/1, 5/1, 6/1, 7/1, 1/2, 8/2, 1/3, 4/3, 5/3, 8/3, 1/4, 8/4, 2/5, 3/5, 4/5, 5/5, 6/5, 7/5}
    \fill [black!80] (\x, -\y) rectangle ++(1, -1);
  \draw [black!25, line width=0.3pt] (0, 0) grid (10, -7);
  \draw [black, line width=0.6pt] (0, 0) rectangle (10, -7);
  \fill [red!80!black] (2.5, -3.5) circle (0.18);
  \node [font=\sffamily\scriptsize, anchor=north] at (5, -7.2)
    {\ifnum#1=1 depois do clique\else antes do clique\fi};
\end{tikzpicture}}

% ------------------------------------------------------------------
% Memória de um programa: objetos e ponteiros. Os objetos que as raízes não
% alcançam são lixo.
\newcommand{\grafoMemoria}{%
\begin{tikzpicture}
  \node [pv, font=\sffamily\small] (raiz) at (0, 1.2) {raízes};
  \node [gv] (a) at (2, 2) {a};
  \node [gv] (b) at (2, 0.4) {b};
  \node [gv] (c) at (3.6, 2) {c};
  \node [gv] (d) at (3.6, 0.4) {d};
  \node [gv, gcinza, dashed] (e) at (5.6, 2) {e};
  \node [gv, gcinza, dashed] (f) at (5.6, 0.4) {f};
  \node [gv, gcinza, dashed] (g) at (7.0, 1.2) {g};
  \draw [arco] (raiz) -- (a);
  \draw [arco] (raiz) -- (b);
  \draw [arco] (a) -- (c);
  \draw [arco] (b) -- (d);
  \draw [arco] (d) -- (c);
  \draw [arco] (e) -- (f);
  \draw [arco] (f) -- (g);
  \draw [arco] (e) -- (c);
  \node [gnota, text=black!60, above=0.1cm of e, xshift=0.7cm] {lixo};
\end{tikzpicture}}


\setminted{
  fontsize=\footnotesize,
  breaklines=true,
  breakanywhere=true,
  frame=lines,
  framesep=2mm,
  baselinestretch=1.05,
  tabsize=2
}

\apptocmd{\frame}{}{\justifying}{}

\setbeamertemplate{navigation symbols}{}

\title{12 - Representação e percursos em grafos}
\author{Prof. Alex Kutzke}
\date{Outubro de 2026}

\begin{document}

% ==========================================
\begin{frame}
	\begin{center}
		\large{Setor de Educação Profissional e Tecnológica}\\
		Tecnologia em Análise e Desenvolvimento de Sistemas\\
		\vspace{1cm}
		\small{DS143 - Estruturas de Dados II}
	\end{center}
	\titlepage
\end{frame}

% ==========================================
\begin{frame}{Objetivos da aula}
	\footnotesize
	Ao final desta aula você deve ser capaz de:
	\begin{enumerate}
		\itemsep=0pt
		\item Usar a terminologia básica de grafos: vértice, aresta, grau,
		      caminho, ciclo, grafo conexo, componente conexa, grafo dirigido;
		\item Representar um grafo por matriz e por listas de adjacência, e
		      escolher entre as duas;
		\item Executar à mão a busca em largura e a busca em profundidade,
		      registrando a ordem de visita e o vetor \texttt{pai};
		\item Explicar por que a busca em largura encontra o caminho com menos
		      arestas, e a busca em profundidade não;
		\item Usar a busca em profundidade para encontrar componentes conexas e
		      detectar ciclos;
		\item Obter uma ordenação topológica pela pós-ordem invertida;
		\item Reconhecer as componentes fortemente conexas de um grafo dirigido
		      pequeno.
	\end{enumerate}
	\vspace{0.1cm}
	Este conteúdo cai na Prova 2.
\end{frame}

% ==========================================
\begin{frame}{Roteiro}
	\tableofcontents
\end{frame}

% ==========================================
\section{Grafos}

\begin{frame}[fragile]{Demonstração: um labirinto}
	\begin{minted}[fontsize=\scriptsize]{text}
$ go run labirinto.go
Busca em largura                 Busca em profundidade
  #####################            #####################
  #S....   #          #            #S.......#          #
  #    .   #   ####   #            #       .#   ####   #
  #####.   #      #   #            #####   .#      #   #
  #  ...   #####  #   #            #  ......#####  #   #
  #  .###             #            #  .###.............#
  #E..  #             #            #E..  #.............#
  #####################            #####################
  13 passos, 43 células visitadas  45 passos, 48 células visitadas
\end{minted}

	\begin{itemize}
		\item a mesma regra nas duas: de uma célula livre se passa para as
		      células livres vizinhas;
		\item a busca em largura achou o caminho mais curto; a em profundidade,
		      um caminho qualquer.
	\end{itemize}
\end{frame}

\begin{frame}{Grafo}
	Um \textbf{grafo} $G = (V, E)$: um conjunto $V$ de \textbf{vértices} e um
	conjunto $E$ de \textbf{arestas}, cada uma ligando dois vértices.

	\begin{center}
		\scalebox{0.85}{\grafoG}
	\end{center}

	\begin{itemize}
		\item 8 vértices, 10 arestas. $V$ e $E$ também indicam as quantidades:
		      ``a busca custa $O(V + E)$'';
		\item o desenho não faz parte do grafo: o que define o grafo é a lista de
		      arestas.
	\end{itemize}
\end{frame}

\begin{frame}{Terminologia}
	\footnotesize
	\begin{tabular}{p{3.2cm}p{4.2cm}p{2.4cm}}
		\toprule
		Termo & Definição & No grafo \\
		\midrule
		adjacentes (vizinhos) & ligados por uma aresta & 2 e 4 \\
		grau & número de arestas que tocam o vértice & 3 tem grau 3 \\
		caminho & sequência de vértices, cada um adjacente ao seguinte & 0, 2, 4, 6 \\
		comprimento & número de arestas do caminho & 3 \\
		ciclo & caminho que volta ao início, sem repetir arestas & 0, 1, 3, 2, 0 \\
		conexo & existe caminho entre quaisquer dois vértices & sim \\
		componente conexa & conjunto máximo de vértices ligados por caminhos & uma só \\
		\bottomrule
	\end{tabular}

	\vspace{0.3cm}
	Sem laço (aresta de $v$ para $v$) e sem arestas paralelas.
\end{frame}

\begin{frame}{Graus, densidade e árvores}
	\begin{itemize}
		\item cada aresta soma 1 ao grau de cada ponta: a soma dos graus é
		      $2E$. No exemplo, $3+2+3+3+3+2+3+1 = 20$;
		\item no máximo $V(V-1)/2$ arestas. Perto disso, grafo \textbf{denso};
		      com $E$ da ordem de $V$, \textbf{esparso};
		\item um cruzamento de ruas tem 4 vizinhos, e não milhares: a maior parte
		      dos grafos de aplicações é esparsa.
	\end{itemize}

	\vspace{0.2cm}
	\begin{itemize}
		\item \textbf{árvore}: grafo conexo e sem ciclos, com $V - 1$ arestas.
		      As árvores das aulas anteriores, com uma raiz escolhida;
		\item em um grafo qualquer pode haver mais de um caminho entre dois
		      vértices. É isso que muda os algoritmos;
		\item \textbf{grafo dirigido}: o arco $v \to w$ vai de $v$ a $w$, e não
		      volta.
	\end{itemize}
\end{frame}

\begin{frame}{Onde aparecem}
	\footnotesize
	\begin{tabular}{lll}
		\toprule
		Aplicação & Vértices & Arestas \\
		\midrule
		malha viária & cruzamentos & ruas; dirigido com mão única \\
		rede de computadores & roteadores & enlaces \\
		circuito eletrônico & portas lógicas & fios \\
		rede social & pessoas & amizade; dirigido se ``seguir'' \\
		jogo de tabuleiro & posições & jogadas permitidas \\
		biologia & proteínas & interação entre proteínas \\
		web & páginas & links, dirigido \\
		artigos científicos & artigos & ``A cita B'', dirigido \\
		dependências & pacotes & ``A importa B'', dirigido \\
		memória de um programa & objetos & ponteiros, dirigido \\
		labirinto, grade & células livres & passagem entre vizinhas \\
		\bottomrule
	\end{tabular}

	\vspace{0.3cm}
	\small
	Exemplos dos slides de Sedgewick e Wayne e desta aula. As duas próximas
	aulas acrescentam um \textbf{peso} a cada aresta. Nesta, as arestas não têm
	peso.
\end{frame}

% ==========================================
\section{Representação}

\begin{frame}[fragile]{Matriz de adjacência}
	\begin{columns}[T]
		\begin{column}{0.6\textwidth}
			\begin{minted}[fontsize=\scriptsize]{go}
type Matriz struct {
	adj [][]bool
}

func (g *Matriz) AdicionaAresta(v, w int) {
	g.adj[v][w] = true
	g.adj[w][v] = true
}
\end{minted}
		\end{column}
		\begin{column}{0.38\textwidth}
			\begin{minted}[fontsize=\scriptsize]{text}
     0 1 2 3 4 5 6 7
  0  . 1 1 . . 1 . .
  1  1 . . 1 . . . .
  2  1 . . 1 1 . . .
  3  . 1 1 . . . 1 .
  4  . . 1 . . 1 1 .
  5  1 . . . 1 . . .
  6  . . . 1 1 . . 1
  7  . . . . . . 1 .
\end{minted}
		\end{column}
	\end{columns}

	\begin{itemize}
		\item vértices numerados de 0 a $V-1$, como os itens do Union-Find;
		\item cada aresta aparece duas vezes: a matriz é simétrica;
		\item a aresta $v$-$w$ existe? Uma consulta, \texttt{adj[v][w]};
		\item $V^2$ posições, quaisquer que sejam as arestas.
	\end{itemize}
\end{frame}

\begin{frame}[fragile]{Listas de adjacência}
	\begin{columns}[T]
		\begin{column}{0.68\textwidth}
			\begin{minted}[fontsize=\scriptsize]{go}
type Grafo struct {
	adj [][]int
}

func NovoGrafo(n int) *Grafo {
	return &Grafo{adj: make([][]int, n)}
}

func (g *Grafo) V() int { return len(g.adj) }

func (g *Grafo) AdicionaAresta(v, w int) {
	g.adj[v] = append(g.adj[v], w)
	g.adj[w] = append(g.adj[w], v)
}
\end{minted}
		\end{column}
		\begin{column}{0.3\textwidth}
			\begin{minted}[fontsize=\scriptsize]{text}
  0: [1 2 5]
  1: [0 3]
  2: [0 3 4]
  3: [1 2 6]
  4: [2 5 6]
  5: [0 4]
  6: [3 4 7]
  7: [6]
\end{minted}
		\end{column}
	\end{columns}

	\begin{itemize}
		\item total de posições nas listas: a soma dos graus, $2E$;
		\item as buscas visitam os vizinhos na ordem da lista.
	\end{itemize}
\end{frame}

\begin{frame}{Matriz ou listas}
	\begin{center}
		\begin{tabular}{lcc}
			\toprule
			Operação & Matriz & Listas \\
			\midrule
			memória & $V^2$ & $V + 2E$ \\
			inserir aresta & $O(1)$ & $O(1)$ \\
			a aresta $v$-$w$ existe? & $O(1)$ & $O(\text{grau}(v))$ \\
			percorrer os vizinhos de $v$ & $O(V)$ & $O(\text{grau}(v))$ \\
			\bottomrule
		\end{tabular}
	\end{center}

	\vspace{0.2cm}
	1 milhão de cruzamentos, 4 vizinhos em média:
	\begin{itemize}
		\item listas: $4 \times 10^6$ inteiros de 8 bytes (32 MB), mais 24 bytes
		      de cabeçalho por \textit{slice} (24 MB);
		\item matriz de \texttt{bool}: $10^{12}$ bytes, 1 TB.
	\end{itemize}

	\vspace{0.2cm}
	Buscas em grafo esparso: listas. Grafo denso, ou muita pergunta ``a aresta
	existe?'': matriz.
\end{frame}

\begin{frame}[fragile]{Lista de arestas e vértices com nome}
	\begin{itemize}
		\item a \textbf{lista de arestas}, como a entrada dos programas, é uma
		      terceira representação. Não serve para buscas; é a entrada do
		      algoritmo de Kruskal, na próxima aula;
		\item vértices com nome: um \texttt{map} de nome para número, e um
		      \textit{slice} para o caminho inverso.
	\end{itemize}

	\begin{minted}[fontsize=\scriptsize]{go}
var indice = map[string]int{}
var nomes []string

func numero(nome string) int {
	if v, ok := indice[nome]; ok {
		return v
	}
	indice[nome] = len(nomes)
	nomes = append(nomes, nome)
	return len(nomes) - 1
}
\end{minted}
\end{frame}

% ==========================================
\section{Buscas}

\begin{frame}[fragile]{Busca em largura (BFS)}
	\begin{minted}[fontsize=\scriptsize]{go}
func BuscaEmLargura(g *Grafo, s int) *BFS {
	b := &BFS{marcado: make([]bool, g.V()),
		pai: make([]int, g.V()), dist: make([]int, g.V())}
	b.pai[s] = -1
	b.marcado[s] = true
	fila := []int{s}
	for len(fila) > 0 {
		v := fila[0]
		fila = fila[1:]
		for _, w := range g.adj[v] {
			if !b.marcado[w] {
				b.marcado[w] = true
				b.pai[w] = v
				b.dist[w] = b.dist[v] + 1
				fila = append(fila, w)
			}
		}
	}
	return b
}
\end{minted}
\end{frame}

\begin{frame}{Busca em largura passo a passo}
	\begin{center}
		\only<1>{\bfsPasso{1}}
		\only<2>{\bfsPasso{2}}
		\only<3>{\bfsPasso{3}}
		\only<4>{\bfsPasso{4}}
		\only<5>{\bfsPasso{5}}
		\only<6>{\bfsPasso{6}}
		\only<7>{\bfsPasso{7}}
		\only<8>{\bfsPasso{8}}
		\only<9>{\bfsPasso{9}}
	\end{center}
\end{frame}

\begin{frame}[fragile]{Busca em largura em $G$, a partir de 0}
	\begin{columns}[T]
		\begin{column}{0.48\textwidth}
			\begin{minted}[fontsize=\scriptsize]{text}
retira 0   fila: [1 2 5]
retira 1   fila: [2 5 3]
retira 2   fila: [5 3 4]
retira 5   fila: [3 4]
retira 3   fila: [4 6]
retira 4   fila: [6]
retira 6   fila: [7]
retira 7   fila: []
\end{minted}
		\end{column}
		\begin{column}{0.48\textwidth}
			\begin{minted}[fontsize=\scriptsize]{text}
  v | dist | pai
  --+------+----
  0 |    0 |  -1
  1 |    1 |   0
  2 |    1 |   0
  3 |    2 |   1
  4 |    2 |   2
  5 |    1 |   0
  6 |    3 |   3
  7 |    4 |   6
\end{minted}
		\end{column}
	\end{columns}

	\begin{itemize}
		\item ao retirar o 2: o 0 e o 3 já estão marcados, só o 4 entra;
		\item ao retirar o 5: nada entra;
		\item é o percurso em largura do material de árvores, com o vetor
		      \texttt{marcado}: em grafo, um vértice pode ser alcançado por mais de
		      um caminho.
	\end{itemize}
\end{frame}

\begin{frame}{Árvore da busca em largura}
	\begin{center}
		\scalebox{0.9}{\grafoGbfs}
	\end{center}

	As arestas \texttt{pai[w]}-\texttt{w} formam uma árvore. Os vértices saem
	da fila em ordem de distância: 0; depois 1, 2, 5; depois 3, 4; depois 6;
	depois 7.

	\vspace{0.2cm}
	Custo: cada vértice entra uma vez na fila, e cada lista é percorrida uma
	vez. $O(V + E)$ com listas, $O(V^2)$ com matriz.
\end{frame}

\begin{frame}[fragile]{Caminho com menos arestas}
	\begin{itemize}
		\item um vértice à distância $d+1$ só entra na fila quando um vizinho à
		      distância $d$ sai, e todos os mais próximos já entraram;
		\item \texttt{dist[v]} é o menor número de arestas de $s$ até $v$.
	\end{itemize}

	\begin{minted}[fontsize=\scriptsize]{go}
func (b *BFS) Caminho(v int) []int {
	if !b.marcado[v] {
		return nil
	}
	var cam []int
	for x := v; x != -1; x = b.pai[x] {
		cam = append(cam, x)
	}
	slices.Reverse(cam)
	return cam
}
\end{minted}

	\begin{itemize}
		\item de 0 até 7: \texttt{[0 1 3 6 7]}, 4 arestas. O caminho 0, 2, 3, 6, 7
		      também tem 4: a busca garante o número de arestas, e não o caminho.
	\end{itemize}
\end{frame}

\begin{frame}{Aplicações da busca em largura}
	\begin{itemize}
		\item \textbf{graus de separação}: conexões de 1\textordmasculine, 2\textordmasculine{} e 3\textordmasculine{} grau do
		      LinkedIn; número de Kevin Bacon (atores ligados por um filme em
		      comum); número de Erd\H{o}s (coautoria de artigos);
		\item \textbf{redes de computadores}: rota com menos saltos
		      (\textit{hops}) entre dois roteadores;
		\item \textbf{coletor da web} (\textit{web crawler}): fila de páginas e
		      conjunto de endereços já vistos; os vizinhos de uma página só
		      aparecem depois que ela é baixada;
		\item \textbf{quebra-cabeça de 8 peças}: cada arrumação do tabuleiro é
		      um vértice, cada movimento é uma aresta; 181.440 arrumações,
		      nenhuma a mais de 31 movimentos da solução.
	\end{itemize}

	\vspace{0.2cm}
	Em todos os casos, cada passo tem o mesmo custo. Passos com custos
	diferentes são o assunto da aula de caminhos mínimos.
\end{frame}

\begin{frame}[fragile]{Busca em profundidade (DFS)}
	\begin{minted}{go}
func (d *DFS) visita(v int) {
	d.marcado[v] = true
	for _, w := range d.g.adj[v] {
		if !d.marcado[w] {
			d.pai[w] = v
			d.visita(w)
		}
	}
}
\end{minted}

	\begin{itemize}
		\item vai o mais longe possível por um caminho, e volta quando não há
		      vizinho novo;
		\item a pilha de chamadas faz o papel da fila da busca em largura;
		\item custo: $O(V + E)$ com listas, pelo mesmo argumento.
	\end{itemize}
\end{frame}

\begin{frame}{Busca em profundidade passo a passo}
	\begin{center}
		\only<1>{\dfsPasso{1}}
		\only<2>{\dfsPasso{2}}
		\only<3>{\dfsPasso{3}}
		\only<4>{\dfsPasso{4}}
		\only<5>{\dfsPasso{5}}
		\only<6>{\dfsPasso{6}}
		\only<7>{\dfsPasso{7}}
		\only<8>{\dfsPasso{8}}
		\only<9>{\dfsPasso{9}}
		\only<10>{\dfsPasso{10}}
		\only<11>{\dfsPasso{11}}
		\only<12>{\dfsPasso{12}}
		\only<13>{\dfsPasso{13}}
		\only<14>{\dfsPasso{14}}
		\only<15>{\dfsPasso{15}}
		\only<16>{\dfsPasso{16}}
	\end{center}
\end{frame}

\begin{frame}[fragile]{Busca em profundidade em $G$, a partir de 0}
	\begin{columns}[T]
		\begin{column}{0.45\textwidth}
			\begin{minted}[fontsize=\scriptsize]{text}
visita(0)
  visita(1)
    visita(3)
      visita(2)
        visita(4)
          visita(5)
          termina 5
          visita(6)
            visita(7)
            termina 7
          termina 6
        termina 4
      termina 2
    termina 3
  termina 1
termina 0
\end{minted}
		\end{column}
		\begin{column}{0.52\textwidth}
			\begin{itemize}
				\small
				\item de 3, lista \texttt{[1 2 6]}: desce para o 2, e não para o 6;
				\item o 6 só é alcançado pelo 4;
				\item \textbf{pré-ordem}, quando a chamada começa:\\
				      \texttt{0 1 3 2 4 5 6 7}
				\item \textbf{pós-ordem}, quando a chamada termina:\\
				      \texttt{5 7 6 4 2 3 1 0}
				\item na pós-ordem, o 4 vem depois de 5, 6 e 7, que ele alcança.
			\end{itemize}
		\end{column}
	\end{columns}
\end{frame}

\begin{frame}{Árvore da busca em profundidade}
	\begin{center}
		\scalebox{0.9}{\grafoGdfs}
	\end{center}

	Caminho de 0 até 7 pelo vetor \texttt{pai}: 0, 1, 3, 2, 4, 6, 7, com 6
	arestas. A busca em largura achou um com 4. A busca em profundidade não dá
	garantia nenhuma sobre o número de arestas.
\end{frame}

% ==========================================
\section{Aplicações}

\begin{frame}[fragile]{Componentes conexas}
	\begin{minted}[fontsize=\scriptsize]{go}
quantas := 0
for v := 0; v < g.V(); v++ {
	if !marcado[v] {
		visita(v, quantas) // marca tudo o que v alcança, com id = quantas
		quantas++
	}
}
\end{minted}

	\begin{center}
		\scalebox{0.8}{\grafoComp}
	\end{center}

	\begin{minted}[fontsize=\scriptsize]{text}
  v  | 0 1 2 3 4 5 6 7 8 9
  id | 0 0 0 1 1 2 2 2 2 3
\end{minted}
\end{frame}

\begin{frame}{Componentes conexas e Union-Find}
	Mesmo papel do \texttt{id} do Quick-Find: $v$ e $w$ estão conectados se, e
	somente se, \texttt{id[v] == id[w]}.

	\vspace{0.3cm}
	\begin{tabular}{p{4.8cm}p{4.8cm}}
		\toprule
		Busca & Union-Find \\
		\midrule
		precisa do grafo inteiro guardado & recebe as arestas uma de cada vez,
		sem guardá-las \\
		$O(V + E)$ uma vez, depois $O(1)$ por consulta & consultas entre uma
		aresta e outra \\
		\bottomrule
	\end{tabular}

	\vspace{0.3cm}
	O algoritmo de Kruskal, na próxima aula, usa o Union-Find desse segundo jeito.
\end{frame}

\begin{frame}{Componentes em imagens}
	\begin{center}
		\imagemBalde{0}\hspace{1cm}\imagemBalde{1}
	\end{center}

	\begin{itemize}
		\item cada pixel é um vértice, ligado aos vizinhos de mesma cor: grafo
		      implícito em grade, como o do labirinto;
		\item \textbf{balde de tinta} (\textit{flood fill}): pinta a componente
		      do pixel clicado; a varinha mágica seleciona da mesma forma;
		\item imagem de microscópio: cada componente de pixels claros é uma
		      partícula, e contar componentes é contar partículas.
	\end{itemize}
\end{frame}

\begin{frame}[fragile]{Detecção de ciclo}
	Na busca em profundidade, a aresta $v$-$w$ examinada leva a três casos:
	\begin{itemize}
		\item $w$ não marcado: desce para $w$;
		\item $w$ marcado e pai de $v$: a aresta pela qual a busca chegou a $v$;
		\item $w$ marcado e não é o pai: outro caminho de $w$ até $v$. A aresta
		      fecha um \textbf{ciclo}.
	\end{itemize}

	\begin{minted}[fontsize=\scriptsize]{text}
$ go run componentes.go
Ciclo: a aresta 2-0 leva a um vértice já marcado
Sem as arestas 1-2 e 7-8, tem ciclo? false
\end{minted}

	A busca desce 0, 1, 2. Na lista do 2, o 0 está marcado e o pai do 2 é o 1:
	ciclo 0, 1, 2.
\end{frame}

% ==========================================
\section{Grafos dirigidos}

\begin{frame}[fragile]{Grafos dirigidos}
	\begin{minted}{go}
func (g *Digrafo) AdicionaArco(v, w int) {
	g.adj[v] = append(g.adj[v], w)
}
\end{minted}

	\begin{itemize}
		\item o arco $v \to w$ entra só na lista de $v$;
		\item as duas buscas funcionam sem mudança, seguindo o sentido dos arcos;
		\item a busca a partir de $s$ marca os vértices \textbf{alcançáveis} a
		      partir de $s$;
		\item $w$ pode ser alcançável a partir de $v$ sem que $v$ seja alcançável
		      a partir de $w$.
	\end{itemize}
\end{frame}

\begin{frame}{Alcançabilidade: coleta de lixo}
	\begin{center}
		\scalebox{0.9}{\grafoMemoria}
	\end{center}

	\begin{itemize}
		\small
		\item objetos são vértices, ponteiros são arcos; as \textbf{raízes} são
		      as variáveis locais e globais;
		\item \textbf{marca e varre} (\textit{mark-sweep}, McCarthy, 1960): marca
		      o que as raízes alcançam e libera o resto;
		\item o coletor de lixo de Go é desse tipo; o código do \texttt{runtime}
		      usa branco, cinza (ponteiros por examinar) e preto;
		\item análise de código: um trecho que o início da função não alcança
		      nunca executa (\texttt{go vet}: \texttt{unreachable code}).
	\end{itemize}
\end{frame}

\begin{frame}{Dependências entre pacotes}
	Arco $B \to A$: o pacote $A$ importa $B$, e $B$ é compilado antes.

	\begin{center}
		\scalebox{0.85}{\grafoPacotes}
	\end{center}

	\textbf{Ordenação topológica}: uma ordem dos vértices em que todo arco
	aponta para a frente. Aqui, uma ordem de compilação.
\end{frame}

\begin{frame}{Pós-ordem invertida}
	\begin{itemize}
		\item grafo sem ciclo: quando a chamada de $v$ termina, tudo o que $v$
		      alcança já terminou, dentro da chamada de $v$ ou antes dela;
		\item na pós-ordem, cada vértice aparece depois de tudo o que ele
		      alcança;
		\item invertida, cada vértice aparece antes de tudo o que ele alcança: é
		      uma ordenação topológica.
	\end{itemize}

	\vspace{0.2cm}
	Uma busca a partir de cada vértice ainda não visitado, como em componentes,
	e no fim a pós-ordem invertida.
\end{frame}

\begin{frame}[fragile]{Ordem de compilação}
	\begin{minted}[fontsize=\scriptsize]{text}
$ go run topologica.go
Ordem de compilação: [log config modelo banco servico http main]
\end{minted}

	\begin{center}
		\scalebox{0.82}{\grafoPacotesOrdem}
	\end{center}

	Todos os arcos apontam para a direita. A ordem não é única:
	\texttt{config log modelo banco servico http main} também serve.
\end{frame}

\begin{frame}{Ciclo em grafo dirigido}
	\texttt{modelo} passa a importar \texttt{servico}:

	\begin{center}
		\scalebox{0.85}{\grafoPacotesCiclo}
	\end{center}

	\begin{itemize}
		\item nenhuma ordem resolve: ordenação topológica existe se, e somente
		      se, não há ciclo (\textbf{DAG}, \textit{directed acyclic graph});
		\item a regra ``marcado e não é o pai'' falha: sem esse arco, a busca
		      a partir de \texttt{modelo} já termina o \texttt{servico} via
		      \texttt{banco}, e depois o arco \texttt{modelo -> servico} chega a
		      um vértice marcado sem ciclo nenhum.
	\end{itemize}
\end{frame}

\begin{frame}[fragile]{Três estados}
	\begin{itemize}
		\item \textbf{branco}: ainda não visitado;
		\item \textbf{cinza}: chamada em andamento, ainda na pilha de chamadas;
		\item \textbf{preto}: chamada terminada.
	\end{itemize}

	\vspace{0.1cm}
	Arco $v \to w$ com $w$ cinza: $v$ foi alcançado a partir de $w$, e o arco
	volta. Ciclo. Arco para vértice preto não fecha ciclo.

	\begin{minted}[fontsize=\scriptsize]{text}
$ go run topologica.go
Com modelo importando servico:
  ciclo: [modelo banco servico]

$ go build .
...
	imports ciclo/b from a.go
	imports ciclo/a from b.go: import cycle not allowed
\end{minted}
\end{frame}

\begin{frame}{Ordenação topológica e ciclos em outras ferramentas}
	\footnotesize
	\begin{tabular}{@{}>{\raggedright\arraybackslash}p{2.0cm}>{\raggedright\arraybackslash}p{1.6cm}>{\raggedright\arraybackslash}p{2.4cm}>{\raggedright\arraybackslash}p{3.5cm}@{}}
		\toprule
		Onde & Vértice & Arco & Ciclo \\
		\midrule
		\texttt{make} & alvo & pré-requisito &
		\texttt{Circular c <- a dependency dropped.} \\
		planilha & célula & fórmula que usa a célula &
		\texttt{Err:522}, referência circular (LibreOffice) \\
		grade curricular & disciplina & pré-requisito &
		nenhum aluno conclui o curso \\
		compilador de Go & tipo & \texttt{struct} dentro de \texttt{struct} &
		\texttt{invalid recursive type} \\
		compilador de Java & classe & herança & \texttt{cyclic inheritance} \\
		\bottomrule
	\end{tabular}

	\vspace{0.3cm}
	\small
	Sem ciclo, uma ordenação topológica dá a ordem de construção, de recálculo
	ou de matrícula.
\end{frame}

\begin{frame}{Componentes fortemente conexas}
	$v$ e $w$ \textbf{fortemente conectados}: cada um alcança o outro.

	\begin{center}
		\scalebox{0.85}{\grafoCFC}
	\end{center}

	\begin{itemize}
		\item do 3 não se volta ao 1; o 6 fica sozinho;
		\item ruas de mão única: em cada componente, vai-se de qualquer
		      cruzamento a qualquer outro e volta-se;
		\item trocando cada componente por um vértice, o grafo não tem ciclo;
		\item Kosaraju-Sharir e Tarjan: $O(V + E)$, com busca em profundidade.
		      Fora desta disciplina.
	\end{itemize}
\end{frame}

% ==========================================
\section{Fechamento}

\begin{frame}[fragile]{De volta ao labirinto: grafo implícito}
	\begin{minted}[fontsize=\scriptsize]{go}
func vizinhos(v int) []int {
	l, c := v/colunas, v%colunas
	var viz []int
	for _, d := range [][2]int{{0, 1}, {1, 0}, {0, -1}, {-1, 0}} {
		nl, nc := l+d[0], c+d[1]
		if mapa[nl][nc] != '#' {
			viz = append(viz, celula(nl, nc))
		}
	}
	return viz
}
\end{minted}

	\begin{itemize}
		\item vértice: a célula \texttt{linha*colunas + coluna};
		\item sem listas guardadas: os vizinhos são calculados quando a busca
		      pede. As buscas trocam \texttt{g.adj[v]} por \texttt{vizinhos(v)};
		\item a grade da etapa 2 do trabalho prático é um grafo desse tipo.
	\end{itemize}
\end{frame}

\begin{frame}{De volta ao labirinto: os dois caminhos}
	\begin{itemize}
		\item busca em largura: 13 passos, porque \texttt{dist[E]} é 13, o menor
		      número de arestas de \texttt{S} até \texttt{E};
		\item busca em profundidade: vizinhos na ordem direita, abaixo,
		      esquerda, acima. Segue pela região aberta da direita antes de voltar
		      para \texttt{E}, e o vetor \texttt{pai} registra 45 passos;
		\item células visitadas: 43 e 48, de 92 livres. Nenhuma das duas
		      garante visitar menos; depende da posição de \texttt{E} e da ordem
		      dos vizinhos.
	\end{itemize}
\end{frame}

\begin{frame}{Escolher entre as buscas}
	\begin{center}
		\begin{tabular}{ll}
			\toprule
			Problema & Busca \\
			\midrule
			caminho com menos arestas & em largura \\
			vértices alcançáveis & qualquer uma \\
			componentes conexas & qualquer uma \\
			ciclo, grafo não dirigido & qualquer uma \\
			ciclo, grafo dirigido & em profundidade, três estados \\
			ordenação topológica & em profundidade, pós-ordem invertida \\
			\bottomrule
		\end{tabular}
	\end{center}
\end{frame}

\begin{frame}{Para a próxima aula}
	\begin{itemize}
		\item exercícios 1 a 11 do roteiro da aula, sem entrega e sem nota. O
		      conteúdo cai na Prova 2;
		\item rodar \texttt{representacao.go}, \texttt{bfs.go}, \texttt{dfs.go},
		      \texttt{componentes.go} e \texttt{topologica.go}, trocando as
		      arestas;
		\item leitura: Sedgewick, seções 4.1 e 4.2; Cormen, capítulo 22.
	\end{itemize}

	\vspace{0.5cm}
	Próxima aula: árvores geradoras mínimas, com os algoritmos de Kruskal e de
	Prim. As arestas passam a ter peso.
\end{frame}

\end{document}
